\chapter{Nasogastric Tube Insertion}

\section*{Overview}
Nasogastric (NG) tube insertion passes a flexible tube through the nose, down the esophagus, and into the stomach for decompression, feeding, or medication administration.

\section*{Alternative Names}
NG tube placement; nasogastric intubation; gastric tube

\section*{Principle}
A lubricated tube is advanced through the nasal passage into the stomach. Correct placement is confirmed by aspiration of gastric contents, pH testing, or radiography. The tube drains the stomach, delivers nutrition, or gives medication.

\section*{History}
Nasogastric tubes have been used since the 19th century, with modern flexible tubes and enteral feeding developed in the 20th century.

\section*{Why It Is Done}
Performed for gastric decompression in bowel obstruction or ileus, short-term enteral feeding, medication administration, gastric lavage after overdose, and obtaining gastric contents for analysis.

\section*{Is This a Primary or Secondary Test or Procedure}
Primary short-term intervention for gastric access or decompression.

\section*{How It Is Performed}
The tube is measured and marked, lubricated, and advanced through the nose while the patient swallows. Placement is verified by aspirating gastric contents, pH testing, or chest radiograph. The tube is secured to the nose.

\section*{Preparation}
Usually none. The patient may be given a local anesthetic spray and sits upright. Coagulation status is reviewed.

\section*{Contraindications and Precautions}
Facial or skull base trauma (risk of intracranial placement), severe coagulopathy, esophageal varices, and recent esophageal surgery are relative contraindications.

\section*{Risks}
Nosebleed, sinusitis, aspiration, esophageal injury, and malposition (including bronchial placement) if position is not confirmed.

\section*{Aftercare and Recovery}
The tube is connected to suction or used for feeds, flushed regularly, and its position is rechecked before each use.

\section*{Results and Interpretation}
Aspiration of acidic gastric fluid and a confirmatory radiograph ensure safe use.

\section*{Expected Results}
A correctly placed tube providing decompression or nutrition without complication.

\section*{Follow-up}
Daily position checks, skin care at the nares, and removal when the indication resolves.

\section*{References}
\begin{enumerate}
  \item MedlinePlus. Nasogastric tube. U.S. National Library of Medicine; 2025.
  \item Current Medical Diagnosis and Treatment. McGraw-Hill; 2025.
\end{enumerate}

\clearpage
